\subsection{Strong-closure theorem contracts: what it means to ``derive'' a continuum field theory}
\label{app:strong_closure_theorem_contracts}

\AuditTag This appendix fixes the theorem-facing meaning of ``derive the continuum Lagrangian/field theory from the finite folding core''.
The goal is to eliminate ambiguity between (i) an \emph{interface representative} closure (a canonical continuum proxy chosen within a finite family) and (ii) a \emph{theorem-level identification} in which a continuum theory is forced---up to explicit equivalences---as the mathematical limit object of a directed finite system.

\subsubsection{Objects, equivalences, and what can be unique}
\paragraph{Action scalar, density, and representative freedom.}
\MathTag A (classical) continuum action is a functional on a field configuration space, typically written as
\[
S[\Phi]=\int \dd^4x\,\mathscr{L}(x,\Phi(x),\partial\Phi(x),\dots),
\qquad
\mathscr{L}=\sqrt{-g}\,\mathcal{L}
\]
in a curved-spacetime dictionary.
Even in purely classical settings, $\mathscr{L}$ is not unique: it is invariant under addition of total divergences and may change under local field redefinitions without altering the Euler--Lagrange equations on-shell.

\begin{definition}[Action equivalence class]
\label{def:action_equivalence_class}
\MathTag Two local Lagrangian densities $\mathscr{L}$ and $\mathscr{L}'$ are said to be \emph{equivalent}, written $\mathscr{L}\sim\mathscr{L}'$, if they differ by:
\begin{itemize}
\item a total divergence $\partial_\mu J^\mu$ of a local current $J^\mu$ (boundary term), and/or
\item a local field redefinition $\Phi\mapsto \Phi'(\Phi)$ that is invertible in the declared configuration class and preserves the locality/regularity assumptions used in the variational calculus.
\end{itemize}
We write $[S]$ for the induced equivalence class of action functionals.
\end{definition}

\paragraph{Canonical representative (when needed).}
\AuditTag When a unique written form is required for auditability, we select a canonical representative in $[S]$ using an explicit finite tie-break rule (e.g.\ the CAP lexicographic complexity key on a declared finite candidate family, as in Appendix~\ref{app:cap_continuum_action_closure}).
Such a choice is a \emph{representative convention} and must not be conflated with theorem-level uniqueness unless a limit theorem discharges the equivalence freedom.

\subsubsection{Identification levels (what ``full physical identification'' means)}
\paragraph{Multiple mathematical packages coexist.}
\AuditTag In this manuscript, three standard packages appear: (i) a protocol-induced quasi-local observable algebra / AQFT local net, (ii) a Wightman field formulation (operator-valued distributions), and (iii) perturbative QFT in the causal Epstein--Glaser / pAQFT carrier.
Each has distinct prerequisites and its own notion of equivalence.

\begin{definition}[Identification levels (contract)]
\label{def:identification_levels_contract}
\MathTag Fix a target continuum theory $\mathcal{T}$ (e.g.\ GR+YM+$\chi$+matter with a specified gauge group and field content).
We say the manuscript achieves:
\begin{itemize}
\item \textbf{(ID-A) Algebraic identification} if the protocol-induced quasi-local algebra/net is isomorphic (in the declared category) to the AQFT net of $\mathcal{T}$, with covariance/microcausality/spectrum properties stated explicitly as assumptions or theorems.
\item \textbf{(ID-W) Wightman identification} if, in addition, there exist generating fields whose vacuum correlators satisfy a Wightman package and reconstruct the same net (Appendices~\ref{app:domain_control_for_generators}--\ref{app:field_reconstruction_theorems}).
\item \textbf{(ID-P) Perturbative identification} if, in addition, the perturbative renormalized time-ordered products and Ward/BRST constraints of $\mathcal{T}$ are realized within a declared truncation and with explicit anomaly filters (Appendices~\ref{app:eg_causal_perturbation_framework}--\ref{app:eg_anomaly_counterterm_restoration} and Appendix~\ref{app:anomaly_theorem_filters}).
\item \textbf{(ID-S) Scattering identification} if, in addition, the scattering theory objects ($S$-matrix / asymptotic states) exist in the stated sense and agree with the continuum dictionary (Appendices~\ref{app:scattering_haag_ruelle_theorems} and \ref{app:scattering_haag_ruelle_lsz_interface}).
\end{itemize}
\end{definition}

\paragraph{Full identification target (this plan).}
\AuditTag The strongest reading ``derive the full continuum field theory'' corresponds to achieving (ID-A)+(ID-W)+(ID-P)+(ID-S), with the understanding that genuinely nonperturbative existence and mass-gap questions remain explicit long-range interface targets (Appendix~\ref{app:nonpert_strong_closure_interface}) unless separately discharged.

\subsubsection{Strong-closure theorem template and failure-point discipline}
\begin{theorem}[Strong-closure target theorem (template)]
\label{thm:strong_closure_target_template}
\MathTag Assume a directed refinement family of finite protocol models is given and satisfies the continuum-limit condition bundle (CL1--CL4) and the algebraic/QFT condition bundles (W1--W3, Spectrum, R1--R3, WBR1--WBR3) as declared in Appendix~\ref{app:minimal_failure_point_templates}.
Then there exists a continuum theory $\mathcal{T}$ and a notion of equivalence $[S]$ (Definition~\ref{def:action_equivalence_class}) such that:
\begin{itemize}
\item the protocol-induced net identifies with the AQFT net of $\mathcal{T}$ (ID-A),
\item whenever the Wightman reconstruction bundle is discharged, the net admits generating fields yielding the Wightman package (ID-W),
\item whenever the EG/pAQFT gauge/BRST and anomaly bundles are discharged to the declared truncation order, the perturbative renormalized dynamics agree with $\mathcal{T}$ (ID-P),
\item whenever the scattering bundles are discharged on the declared window, the scattering objects agree (ID-S).
\end{itemize}
\end{theorem}

\AuditTag Theorem~\ref{thm:strong_closure_target_template} is the contract-level statement: it is only as strong as the explicitly discharged condition bundles.
Whenever any failure point triggers, the manuscript must downgrade to the corresponding fallback recorded in Appendix~\ref{app:minimal_failure_point_templates}.

